% Kiến trúc hệ thống, triển khai từ đầu, bàn giao — cho người kỹ thuật của Pi.
% Dựng: cd docs/guide && xelatex -interaction=nonstopmode -halt-on-error kientruc.tex (hai lần)
\documentclass[aspectratio=169,11pt]{beamer}
\usepackage{fontspec}
\setsansfont{TeX Gyre Heros}
\setmainfont{TeX Gyre Heros}
\newfontfamily\symf{DejaVu Sans}
\usepackage{tikz}
\usetikzlibrary{decorations.pathreplacing,shapes.geometric,shapes.symbols,positioning,arrows.meta,calc,fit,backgrounds,angles,quotes,shadows}
\usepackage{xcolor}

% ---------- Colours ----------
\definecolor{TBT}{HTML}{7A1F5C}
\definecolor{PT}{HTML}{B8541B}
\definecolor{NCB}{HTML}{1F6FB2}
\definecolor{PB}{HTML}{2E8B57}
\definecolor{VP}{HTML}{A67C00}
\definecolor{BTK}{HTML}{5B4FA8}
\definecolor{TG}{HTML}{8B5A2B}
\definecolor{HT}{HTML}{8A8F98}
\definecolor{ink}{HTML}{1E2A38}
\definecolor{soft}{HTML}{F4F6F8}
\definecolor{warn}{HTML}{C0392B}
\definecolor{TAGG}{HTML}{6B7280}

% ---------- Beamer look ----------
\setbeamertemplate{navigation symbols}{}
\setbeamercolor{normal text}{fg=ink}
\setbeamercolor{frametitle}{fg=ink}
\setbeamerfont{frametitle}{size=\large,series=\bfseries}
\setbeamertemplate{frametitle}{%
  \vskip6pt\hskip-2pt\insertframetitle\par
  \vskip2pt{\color{ink!25}\rule{\linewidth}{0.6pt}}\vskip-2pt}
\setbeamertemplate{footline}{%
  \hskip8pt{\tiny\color{ink!45}Đề ra kỳ này · Hướng dẫn sử dụng · bản thử kỳ 10/2026}\hfill
  {\tiny\color{ink!45}\insertframenumber}\hskip8pt\vskip5pt}
\setbeamersize{text margin left=8mm,text margin right=8mm}

% ---------- Icons ----------
\newcommand\human[1][ink]{\tikz[baseline=-0.55ex,x=1cm,y=1cm,sharp corners,solid,line width=0.4pt,scale=0.11]{\fill[#1] (0,0.9) circle (0.55);\fill[#1] (-0.95,-0.9) to[out=90,in=180] (0,0.15) to[out=0,in=90] (0.95,-0.9) -- cycle;}}
\newcommand\gear{{\symf ⚙}}
\newcommand\yes{{\symf ✓}}
\newcommand\no{{\symf ✕}}
\newcommand\fullc{{\symf ●}}
\newcommand\emptyc{{\symf ○}}
\newcommand\lockic{\tikz[baseline=-0.4ex,x=1cm,y=1cm,sharp corners,solid,line width=0.4pt,scale=0.12]{\draw[line width=0.7pt] (-0.55,0.3) -- (-0.55,0.75) arc (180:0:0.55) -- (0.55,0.3);\fill (-0.85,-0.75) rectangle (0.85,0.35);}}
% channel icons (draw at given size)
\newcommand\icmail[1][ink]{\tikz[baseline=-0.6ex,x=1cm,y=1cm,sharp corners,solid,line width=0.4pt,scale=0.16]{\draw[#1,line width=0.6pt,fill=white] (-1,-0.65) rectangle (1,0.65);\draw[#1,line width=0.6pt] (-1,0.65) -- (0,-0.1) -- (1,0.65);}}
\newcommand\icfolder[1][ink]{\tikz[baseline=-0.6ex,x=1cm,y=1cm,sharp corners,solid,line width=0.4pt,scale=0.16]{\draw[#1,line width=0.6pt,fill=#1!15] (-1,-0.7) -- (-1,0.7) -- (-0.35,0.7) -- (-0.15,0.45) -- (1,0.45) -- (1,-0.7) -- cycle;}}
\newcommand\icdoc[2][ink]{\tikz[baseline=-0.6ex,x=1cm,y=1cm,sharp corners,solid,line width=0.4pt,scale=0.16]{\draw[#1,line width=0.6pt,fill=white] (-0.75,-0.95) -- (-0.75,0.95) -- (0.35,0.95) -- (0.75,0.55) -- (0.75,-0.95) -- cycle;\node[font=\fontsize{3.2}{3}\selectfont\bfseries,text=#1] at (0,-0.2) {#2};}}
\newcommand\icweb[1][ink]{\tikz[baseline=-0.6ex,x=1cm,y=1cm,sharp corners,solid,line width=0.4pt,scale=0.16]{\draw[#1,line width=0.6pt,fill=white] (-1.1,-0.75) rectangle (1.1,0.75);\fill[#1] (-1.1,0.42) rectangle (1.1,0.75);\fill[white] (-0.95,0.58) circle (0.08) (-0.75,0.58) circle (0.08);}}
\newcommand\icoverleaf{\icdoc[NCB]{TeX}}
\newcommand\icpdf{\icdoc[warn]{PDF}}

% ---------- Diagram styles ----------
\tikzset{
  >=Stealth,
  hstep/.style={draw=#1,fill=#1!8,line width=0.9pt,rounded corners=2.5pt,align=center,font=\scriptsize,inner sep=3pt,text=ink},
  sstep/.style={draw=HT,fill=HT!9,dashed,line width=0.6pt,rounded corners=2.5pt,align=center,font=\tiny,inner sep=2pt,text=HT!55!black},
  dec/.style={diamond,aspect=1.8,draw=#1,fill=#1!8,line width=0.9pt,align=center,font=\scriptsize,inner sep=0.5pt,text=ink},
  tag/.style={fill=#1,text=white,font=\fontsize{5.5}{6}\selectfont\bfseries,rounded corners=1.5pt,inner sep=1.6pt},
  note/.style={draw=ink!25,fill=soft,rounded corners=2pt,font=\tiny,align=left,text=ink!80,inner sep=3pt},
  flow/.style={->,line width=0.7pt,ink!70},
  doc/.style={->,dashed,line width=0.6pt,ink!50},
  lane/.style={fill=#1!5,draw=none},
  lanelab/.style 2 args={fill=#1,text=white,font=\fontsize{6}{7}\selectfont\bfseries,minimum width=9mm,minimum height=#2,inner sep=0pt,align=center},
  badge/.style={circle,fill=#1,text=white,font=\scriptsize\bfseries,minimum size=8mm,inner sep=0pt},
  mtag/.style={draw=TAGG,fill=TAGG!10,text=TAGG!70!black,font=\fontsize{6}{7}\selectfont\ttfamily,rounded corners=1.5pt,inner sep=1.8pt},
  stag/.style={fill=PT,text=white,font=\fontsize{6.5}{7}\selectfont\bfseries,rounded corners=1.5pt,inner sep=2pt},
}
% owner tag on a node's top-right corner
\newcommand\owner[2]{\node[tag=#1,anchor=center] at ([xshift=-2pt]#2.north east) {#1};}
\newcommand\ownerHT[1]{\node[tag=HT,anchor=center] at ([xshift=-2pt]#1.north east) {HT};}
% a decision with PT on top, TBT approval below
\newcommand\decTags[1]{\node[tag=PT,anchor=center] at (#1.north) {PT};\node[tag=TBT,anchor=center] at (#1.south) {TBT duyệt};}

% "you are here" strip, top-right of frame
\newcommand\youare[1]{%
\begin{tikzpicture}[remember picture,overlay]
\node[anchor=north east,inner sep=0pt] at ([xshift=-6mm,yshift=-3.2mm]current page.north east){%
\begin{tikzpicture}[x=6.2mm,font=\fontsize{4.5}{5}\selectfont]
\foreach \i/\l/\c in {1/Gửi bài/TG,2/LaTeX/NCB,3/Sơ bộ/NCB,4/Shortlist/NCB,5/Phản biện/PB,6/Chọn/PT,7/Khoá/PT,8/In/BTK}{
  \ifnum\i=#1\relax
    \node[circle,fill=\c,minimum size=2.6mm,inner sep=0] (y\i) at (\i,0) {};
    \node[below=0.2mm of y\i,text=\c,font=\fontsize{4.5}{5}\selectfont\bfseries] {\l};
  \else
    \node[circle,fill=\c!35,minimum size=1.7mm,inner sep=0] (y\i) at (\i,0) {};
    \node[below=0.3mm of y\i,text=ink!45] {\l};
  \fi}
\foreach \i [count=\j from 2] in {1,...,7}{\draw[ink!25,line width=0.4pt] (y\i) -- (y\j);}
\end{tikzpicture}};
\end{tikzpicture}}

% ---------- P1036 figure (verified configuration) ----------
\newcommand\figP[2][1]{%
\begin{tikzpicture}[x=1cm,y=1cm,scale=#1,line cap=round,line join=round,sharp corners,solid]
\coordinate (A) at (-0.4243,1.3343);\coordinate (B) at (-0.7660,-0.6428);
\coordinate (C) at (0.1736,-0.9848);\coordinate (D) at (0.7660,0.6428);
\coordinate (J) at (-0.3241,0.9460);\coordinate (T) at (2.1260,1.7843);
\draw[ink!35,line width=0.6pt] (0,0) circle (1);
\draw[NCB,line width=0.9pt] (A)--(B)--(C)--(D)--cycle;
\draw[ink!40,line width=0.5pt] (A)--(C) (B)--(T);
\draw[PT!80,line width=0.6pt] ($(J)!-0.22!(T)$)--(T);
\draw[warn,line width=1.2pt] (T)--(J) (T)--(A);
\path (T)--(J) node[pos=0.5,sloped,warn,font=\tiny]{$|$} (T)--(A) node[pos=0.5,sloped,warn,font=\tiny]{$|$};
\pic[draw=ink!70,angle radius=3.2mm,"$\scriptscriptstyle\alpha$",angle eccentricity=1.6]{angle=C--B--A};
\pic[draw=ink!70,angle radius=3.2mm,"$\scriptscriptstyle\alpha$",angle eccentricity=1.6]{angle=A--D--C};
\foreach \p/\pos in {A/above left,B/below left,C/below,D/right,T/above right}\fill (\p) circle (0.9pt) node[\pos,font=#2,inner sep=1pt]{$\p$};
\fill (J) circle (0.9pt) node[below left,font=#2,inner sep=0.5pt]{$J$};
\end{tikzpicture}}

\newcommand\sumbox[7]{% x,y,w,h,colour,title,body
  \draw[#5,line width=0.8pt,fill=#5!6,rounded corners=3pt] (#1,#2) rectangle ++(#3,-#4);
  \node[font=\scriptsize\bfseries,anchor=north west,text=#5!85!black] at (#1+2,#2-1.5) {#6};
  \node[font=\fontsize{6.6}{8}\selectfont,anchor=north west,align=left,text width=#3*1mm-4mm] at (#1+2,#2-6) {#7};}

\tikzset{cs/.style={hstep=#1,text width=25mm,minimum height=7mm,font=\fontsize{6.3}{7.3}\selectfont}}

% ---------- extra styles for slides 6-15 ----------
\tikzset{
  big/.style={hstep=#1,line width=1.4pt,font=\fontsize{6.6}{7.6}\selectfont\bfseries},
  cyl/.style={cylinder,shape border rotate=90,aspect=0.25,draw=#1,fill=#1!8,line width=0.8pt,align=center,font=\fontsize{6}{7}\selectfont,minimum width=14mm,minimum height=10mm},
  chip/.style={draw=#1,fill=#1!10,rounded corners=1.5pt,font=\fontsize{5.6}{6.5}\selectfont,inner sep=1.6pt,text=ink},
  card/.style={draw=ink!25,fill=white,rounded corners=1.5pt,minimum width=9.6mm,minimum height=9.5mm,inner sep=0.8pt,font=\fontsize{4.3}{5.1}\selectfont,align=center},
  slot/.style={draw=ink!20,dashed,fill=soft,rounded corners=1.5pt,minimum width=10.6mm,minimum height=25mm},
  win/.style={draw=ink!40,fill=white,line width=0.6pt,rounded corners=2pt},
  field/.style={draw=ink!25,fill=white,rounded corners=1pt,font=\fontsize{5.5}{6.5}\selectfont,inner sep=1.5pt,anchor=north west},
  bar/.style={fill=#1,draw=none,rounded corners=1pt},
}
% topic colour stripes
\definecolor{tDS}{HTML}{3B7DD8}\definecolor{tSH}{HTML}{D8A33B}\definecolor{tHH}{HTML}{3BA87A}\definecolor{tTH}{HTML}{B24C8F}
% a selection-board card: #1 pos, #2 topic colour, #3 text
\newcommand\bcard[3]{\node[card] (cc) at #1 {#3};\fill[#2] (cc.north west) rectangle ([yshift=-1mm]cc.north east);}
\newcommand\cellm[1]{{\symf #1}}

\def\GX{36}\def\GW{25}
\newcommand\ganttrow[5]{% y, label, colour, x0(weeks), x1(weeks)
  \node[anchor=east,font=\fontsize{5.8}{6.8}\selectfont,text=#3!80!black] at (\GX-1,#1) {#2};
  \fill[#3,rounded corners=1pt] (\GX+#4*\GW,#1+1.4) rectangle (\GX+#5*\GW,#1-1.4);}

\newcommand\cmt[4]{\node[draw=#2!50,fill=#2!5,rounded corners=1.5pt,anchor=north west,text width=58mm,font=\fontsize{5.4}{6.6}\selectfont,inner sep=2pt] at #1 {\textbf{\textcolor{#2!80!black}{#3}}\\#4};}


\colorlet{PBL}{PB!60}
\newcommand\cn[1]{\tikz[baseline=-0.6ex,x=1cm,y=1cm,sharp corners,solid]\node[circle,draw=ink!70,line width=0.4pt,inner sep=0.3pt,minimum size=2.6mm,font=\fontsize{4.6}{4.6}\selectfont\bfseries]{#1};}

\tikzset{ps/.style={hstep=#1,font=\fontsize{5.8}{6.8}\selectfont,minimum height=6.8mm}}

\definecolor{aChon}{HTML}{00897B}\definecolor{aLoai}{HTML}{C62828}\definecolor{aSom}{HTML}{D81B60}\definecolor{aDuTru}{HTML}{8E24AA}

\tikzset{rl/.style={hstep=#1,font=\fontsize{5.6}{6.6}\selectfont,text width=27mm,align=left,anchor=west,minimum height=5.6mm}}

\newcommand\dpanel[4]{% x, number, title, body
  \draw[ink!25,fill=white,rounded corners=2.5pt] (#1,0) rectangle ++(34,-60);
  \node[circle,fill=ink,text=white,font=\scriptsize\bfseries,minimum size=6mm,inner sep=0] at (#1+5,-5.5) {#2};
  \node[anchor=west,font=\scriptsize\bfseries,text width=24mm,align=left] at (#1+9,-5.5) {#3};
  \node[anchor=north west,font=\fontsize{6}{7.4}\selectfont,text width=31mm,align=left] at (#1+1.5,-39) {#4};}


\tikzset{seg/.style={draw=#1,fill=#1!10,line width=0.8pt,font=\small\ttfamily\bfseries,inner sep=2pt,minimum height=7mm,anchor=west,text=ink}}

% ---------- guide extras ----------
\definecolor{QT}{HTML}{37474F}
\newcommand\fsz[2]{\fontsize{#1}{#2}\selectfont}
\newcommand\rtag[1]{\tikz[baseline=-0.6ex]\node[tag=#1,font=\fsz{8}{9}\bfseries,inner sep=2.2pt]{#1};}
\newcommand\soon{\tikz[baseline=-0.6ex]\node[draw=warn,text=warn,rounded corners=1pt,font=\fsz{5}{6}\bfseries,inner sep=1.2pt]{SẮP CÓ};}
\newcommand\kbd[1]{\tikz[baseline=-0.6ex]\node[draw=ink!40,fill=soft,rounded corners=1pt,font=\fsz{5.6}{6.4}\ttfamily,inner sep=1.3pt]{#1};}
\tikzset{
  btn/.style={fill=NCB,text=white,rounded corners=1.5pt,font=\fsz{5.6}{6.4}\bfseries,inner sep=2pt},
  sel/.style={draw=ink!35,fill=white,rounded corners=1pt,font=\fsz{4.8}{5.6},inner sep=1.2pt},
  pchip/.style={fill=ink!8,rounded corners=1.5pt,font=\fsz{4.8}{5.6},inner sep=1.2pt},
  wchip/.style={fill=warn!12,text=warn!80!black,rounded corners=1.5pt,font=\fsz{4.8}{5.6},inner sep=1.2pt},
  callout/.style={draw=#1,fill=#1!6,rounded corners=1.5pt,font=\fsz{5.6}{6.8},align=left,inner sep=2pt,text=ink},
  cline/.style={#1,line width=0.6pt,-{Stealth[length=1.6mm]}},
  gbar/.style={fill=NCB!10,draw=none},
}
% a browser window with title bar: #1 lower-left x, #2 top y, #3 width, #4 height
\newcommand\browser[4]{%
  \draw[win] (#1,#2) rectangle ++(#3,-#4);
  \fill[ink!8] (#1,#2) rectangle ++(#3,-3.2);
  \fill[warn!70] (#1+2,#2-1.6) circle (0.55); \fill[VP!80] (#1+3.8,#2-1.6) circle (0.55); \fill[PB!80] (#1+5.6,#2-1.6) circle (0.55);
  \node[draw=ink!20,fill=white,rounded corners=1pt,font=\fsz{3.8}{4},inner sep=0.8pt,minimum width=#3*0.4mm] at (#1+#3/2,#2-1.6) {script.google.com};}

\hypersetup{pdftitle={Đề ra kỳ này — Kiến trúc, triển khai, bàn giao},pdfauthor={}}
\setbeamertemplate{footline}{%
  \hskip8pt{\tiny\color{ink!45}Đề ra kỳ này · Kiến trúc, triển khai, bàn giao · 10/2026}\hfill
  {\tiny\color{ink!45}\insertframenumber}\hskip8pt\vskip5pt}
\tikzset{
  box/.style={draw=#1,fill=#1!7,line width=0.8pt,rounded corners=2.5pt,align=center,font=\fontsize{6.4}{7.6}\selectfont,inner sep=2.5pt,text=ink},
  zone/.style={draw=#1!55,dashed,line width=0.6pt,rounded corners=4pt,fill=#1!3},
  zlab/.style={font=\fontsize{6.2}{7}\selectfont\bfseries,text=#1!80!black,anchor=north west},
  ar/.style={->,line width=0.7pt,ink!65},
  al/.style={font=\fontsize{5.2}{6}\selectfont,text=ink!75,fill=white,inner sep=1pt},
}
\colorlet{GG}{NCB}\colorlet{GH}{BTK}\colorlet{MAC}{PT}
\newcommand\sm[1]{{\fontsize{6.6}{8.2}\selectfont #1}}
\def\msg#1#2#3#4{\draw[ar] (#1,#3) -- node[al,above]{#4} (#2,#3);}
\def\grp#1#2#3#4#5{\node[box=#3,anchor=north west,text width=#4mm,align=left] at (#1,#2) {#5};}
\def\crow#1#2#3{\node[anchor=west,font=\fontsize{6.2}{7}\selectfont\ttfamily\bfseries,text=ink] at (0,#1) {#2};
               \node[anchor=west,font=\fontsize{6.2}{7.4}\selectfont,text width=105mm] at (34,#1) {#3};}
\def\ccol#1#2#3#4{\node[box=#2,anchor=north west,text width=44mm,align=left] at (#1,0) {\textbf{#3}\\[1mm] #4};}
\begin{document}

%======================================================== 0
\begin{frame}[plain]
\vfill
{\LARGE\bfseries Đề ra kỳ này}\\[2mm]
{\large Hệ thống được xây thế nào, nằm ở đâu, chạy ra sao}\\[1mm]
{\normalsize\color{ink!65} Kiến trúc · triển khai từ đầu · nắm quyền mã nguồn và dữ liệu}\\[8mm]
\sm{Tài liệu chi tiết trong kho mã: \texttt{docs/setup.md} (cài đặt), \texttt{docs/data-model.md} (dữ liệu),
\texttt{docs/data-policy.md} (quy định dữ liệu), \texttt{docs/testing.md} (kiểm thử).}
\vfill
\end{frame}

%======================================================== 1. Toàn cảnh
\begin{frame}{Toàn cảnh}
\centering
\begin{tikzpicture}[x=1mm,y=1mm]
% vùng Google
\draw[zone=GG] (24,2) rectangle (112,-60); \node[zlab=GG] at (25,1) {Google (tài khoản chủ của Pi)};
\draw[zone=GH] (116,2) rectangle (146,-60); \node[zlab=GH] at (117,1) {GitHub};
\draw[zone=MAC] (-4,-36) rectangle (21,-60); \node[zlab=MAC] at (-3,-37) {Máy quản trị};
% người dùng
\node[box=HT,minimum width=22mm,minimum height=16mm] (u) at (8,-12) {\human\ \human\ \human\\[1mm] Ban biên tập,\\ phản biện\\ \textcolor{ink!60}{trình duyệt, điện thoại}};
% hai ứng dụng
\node[box=GG,text width=24mm] (a) at (40,-10) {\textbf{Ứng dụng đăng nhập}\\ Apps Script \texttt{apps/signin}\\ \textcolor{ink!60}{chỉ hỏi email, ký liên kết}};
\node[box=GG,text width=30mm,minimum height=20mm] (b) at (78,-12) {\textbf{Ứng dụng chính}\\ Apps Script \texttt{apps/main}\\ trang web (\texttt{ui/}) + API\\ \textcolor{ink!60}{quyền theo vai trò ở mỗi lời gọi}};
% dữ liệu
\node[cyl=GG] (s) at (40,-40) {Google Sheet\\ \textbf{dữ liệu}\\ \textcolor{ink!60}{17 tab}};
\node[box=GG,text width=22mm] (d) at (64,-42) {Google Drive\\ thư mục \textbf{hình}, \textbf{chế bản},\\ tệp nhập / vá};
\node[box=GG,text width=20mm] (m) at (98,-36) {Gmail (MailApp)\\ thư mời, nhắc hạn,\\ kết quả kiểm thử};
\node[box=GG,text width=20mm] (t) at (98,-51) {Trigger hẹn giờ\\ \textcolor{ink!60}{nhắc hạn · kiểm thử · đồng bộ hình}};
% github
\node[box=GH,text width=26mm] (r1) at (131,-12) {\textbf{nghia71/pi-drkn}\\ \textcolor{ink!60}{công khai}\\ mã, tài liệu, công cụ,\\ kiểm thử tự động (CI)};
\node[box=GH,text width=26mm] (r2) at (131,-42) {\textbf{pi-drkn-hinh}\\ \textcolor{ink!60}{riêng tư}\\ Actions: TikZ → SVG};
% máy
\node[box=MAC,text width=21mm,font=\fontsize{5.8}{7}\selectfont] (mac) at (8.5,-50) {\texttt{git pull}, \texttt{deploy.sh}\\ gói chế bản (\texttt{khoaky})\\ dựng hình tay};
% mũi tên
\draw[ar] (u) -- node[al,above]{1. mở} (a);
\draw[ar] (a) -- node[al,above]{2. liên kết ký} (b);
\draw[ar,<->] (u.south east) to[out=-20,in=200] node[al,pos=0.55]{3. làm việc} (b.south west);
\draw[ar,<->] (b.south) ++(-8,0) -- (s.north east);
\draw[ar,<->] (b.south) -- (d.north);
\draw[ar] (b.south east) -- (m.north);
\draw[ar] (t.west) to[out=180,in=-60] (b.south);
\draw[ar,<->] (b.east) ++(0,-6) -- node[al,below]{mã TikZ / SVG} (r2.north west);
\draw[ar,dashed] (r1.west) -- node[al,above]{mã} (b.east);
\draw[ar,dashed] (mac.north) to[out=80,in=-120] node[al,pos=0.45]{clasp push} (a.south west);
\draw[ar,dashed] (r1.south) -- node[al,right]{công cụ dựng} (r2.north);
\end{tikzpicture}

\vspace{-1mm}\sm{Không có máy chủ, cơ sở dữ liệu riêng: hai dự án Apps Script chạy bằng quyền của tài khoản chủ; dữ liệu là một Google Sheet
và hai thư mục Drive. Nét đứt: mã đi từ GitHub qua máy quản trị (\texttt{git pull}, \texttt{scripts/deploy.sh}) vào Apps Script.}
\end{frame}

%======================================================== 2. Đăng nhập, quyền
\begin{frame}{Đăng nhập và quyền: Google xác nhận email, hệ thống quyết định vai trò}
\begin{tikzpicture}[x=1mm,y=0.86mm]
\foreach \x/\n/\c in {0/Người dùng/HT,38/Ứng dụng đăng nhập/GG,78/Ứng dụng chính/GG,118/Sheet: tab Users/GG}{
  \node[box=\c,minimum width=26mm] at (\x+12,0) {\textbf{\n}};
  \draw[ink!25,line width=0.6pt] (\x+12,-3) -- (\x+12,-56);}
\msg{12}{50}{-9}{mở địa chỉ /exec (đăng nhập Google)}
\node[note,text width=32mm,anchor=west] at (52,-15) {Google chỉ cho biết \textbf{email}. Ký $s=\mathrm{HMAC\text{-}SHA256}(\text{email}|t,\ \texttt{SECRET})$};
\msg{50}{12}{-22}{trang „Xin chào …" + nút Vào hệ thống}
\msg{12}{90}{-29}{liên kết \texttt{?u=email\&t=…\&s=…}}
\node[note,text width=30mm,anchor=west] at (92,-35) {Kiểm tra chữ ký (cùng \texttt{SECRET}), $t$ không quá 10 phút → phiên ngẫu nhiên, 6 giờ};
\msg{90}{130}{-42}{vai trò? đang hoạt động?}
\msg{90}{12}{-49}{trang web + phiên}
\msg{12}{90}{-55}{mỗi thao tác: \texttt{api(phiên, việc, …)} → kiểm tra vai trò lại}
\end{tikzpicture}

\sm{\textbf{Vì sao hai ứng dụng}: ứng dụng chính chạy bằng quyền của chủ nên không tự biết người xem là ai; ứng dụng đăng nhập chỉ có quyền
xem email. \textbf{Quyền nằm ở máy chủ}: trang web chỉ hiện nút, mọi lời gọi đều kiểm tra lại vai trò (tab Users, có hiệu lực ngay);
phản biện chỉ nhận dữ liệu bài được giao, ẩn tên tác giả. Mọi thao tác ghi vào tab \textbf{Audit}.}
\end{frame}

%======================================================== 3. Dữ liệu
\begin{frame}{Dữ liệu nằm ở đâu}
\begin{tikzpicture}[x=1mm,y=1mm]
\grp{0}{0}{NCB}{44}{\textbf{Bài} (Sheet)\\ \texttt{Problems} — đề, lời giải (bản biên tập và bản gốc), hình, trạng thái, phiên bản\\ \texttt{Authors} — tên in, đơn vị, \textbf{liên hệ} (hạn chế)\\ \texttt{Provenance}, \texttt{ConversionLog} — nguồn, chuyển đổi\\ \texttt{Revisions} — mọi lần sửa\\ \texttt{Corrections}, \texttt{Checks}, \texttt{Conflicts}}
\grp{48}{0}{PB}{42}{\textbf{Phản biện} (Sheet)\\ \texttt{Rounds} — kỳ, hạn, mở lời giải\\ \texttt{Assignments} — giao bài, mời, nhắc\\ \texttt{Reviews} — phiếu phản biện\\ \texttt{Comments} — thảo luận}
\grp{94}{0}{PT}{46}{\textbf{Chọn và in} (Sheet)\\ \texttt{Issues}, \texttt{Shortlist} — bảng chọn bài\\ \texttt{Published} — số in đã dùng\\[1mm] \textbf{Hệ thống} (Sheet)\\ \texttt{Users} — người dùng, vai trò\\ \texttt{Audit} — nhật ký}
\grp{0}{-33}{GG}{66}{\textbf{Drive} (tài khoản chủ, không chia sẻ)\\ \textbf{Pi ĐRKN — hình}: ảnh của bài, \texttt{tikz-<mã>.svg}\\ \textbf{Pi ĐRKN — chế bản}: gói \texttt{.tex + pic/} của mỗi số đã khoá\\ tệp nhập / vá \texttt{pi-drkn-import…}, \texttt{pi-drkn-patch…}}
\grp{70}{-33}{HT}{70}{\textbf{Script properties} (cấu hình, bí mật — không ở trong mã)\\ \texttt{SHEET\_ID} \texttt{SECRET} \texttt{SIGNIN\_URL} \texttt{FIG\_FOLDER\_ID} \texttt{EXPORT\_FOLDER\_ID}\\ \texttt{GITHUB\_REPO} \texttt{GITHUB\_TOKEN} \texttt{TEST\_SHEET\_ID} \texttt{TEST\_USERS}\\ tuỳ chọn: \texttt{BLIND\_REVIEW} \texttt{REMINDER\_DAYS/HOUR} \texttt{BOARD\_LAYOUT} \texttt{TEST\_HOUR}\\ dự án đăng nhập: \texttt{SECRET}, \texttt{B\_URL}}
\end{tikzpicture}

\sm{\textbf{Ra khỏi Google chỉ có}: mã TikZ của hình (kho GitHub riêng tư, để dựng SVG) và thư gửi người dùng (không có đề, tên tác giả).
Kho mã công khai \textbf{không bao giờ} chứa đề, lời giải, tên, email thật: \texttt{scripts/guard.py} chặn ở mỗi commit và trên GitHub.
Kiểm thử dùng Sheet riêng „Pi ĐRKN — kiểm thử" với dữ liệu bịa.}
\end{frame}

%======================================================== 4. Mã nguồn
\begin{frame}{Mã nguồn: kho \texttt{nghia71/pi-drkn}}
\begin{tikzpicture}[x=1mm,y=1mm]
\crow{0}{apps/signin/}{ứng dụng đăng nhập (một tệp \texttt{Code.gs})}
\crow{-6.5}{apps/main/*.gs}{máy chủ: \texttt{Web} (vào hệ thống) · \texttt{Auth} (phiên, vai trò) · \texttt{Api} (mọi thao tác của trang) · \texttt{Db}, \texttt{Schema} (Sheet) ·
 \texttt{Rounds} (kỳ phản biện, thư) · \texttt{Board} (bảng chọn bài) · \texttt{Close} (khoá kỳ, tệp .tex) · \texttt{Fig}, \texttt{FigCloud} (hình) ·
 \texttt{Intake} (thêm bài) · \texttt{Import}, \texttt{Patch} (nhập, vá) · \texttt{Setup}, \texttt{Backup} (sao lưu) · \texttt{Tests} (111 kịch bản)}
\crow{-14}{apps/main/ui/}{trang web: \texttt{Index.html} (một trang), \texttt{Styles.html}, \texttt{Icons.html} (Lucide)}
\crow{-19}{shared/pi-render.js}{hiển thị LaTeX của Pi trên trang (macro, công thức qua MathJax, chặn HTML)}
\crow{-24}{templates/}{\texttt{dinhdang.tex} (mẫu in của Pi), \texttt{hinh-mau.tex} (mẫu dựng hình)}
\crow{-29.5}{tools/}{\texttt{khoaky} (gói chế bản → PDF) · \texttt{hinh} (TikZ → SVG) · \texttt{convert} (chuyển hồ sơ cũ) · \texttt{gas-sim} (Apps Script mô phỏng)}
\crow{-36}{tests/}{bộ hiển thị · gói chế bản · hình · \textbf{giao diện trong Chromium} (\texttt{ui/run.js}) · từ bên ngoài; \texttt{fixtures/}: bài bịa}
\crow{-41.5}{scripts/}{\texttt{deploy.sh} (kiểm thử rồi đẩy) · \texttt{test-all.sh} · \texttt{guard.py} (chặn dữ liệu)}
\crow{-46.5}{docs/}{cài đặt, dữ liệu, quy định, kiểm thử; \texttt{guide/} — các bộ slide này (Beamer)}
\crow{-51.5}{.github/workflows/}{trên GitHub mỗi lần đẩy: chặn dữ liệu, 110 kịch bản máy chủ, kiểm thử giao diện}
\end{tikzpicture}

\sm{Không có thư viện ngoài phía máy chủ; trang web chỉ tải MathJax. Chú thích, giao diện, tài liệu bằng tiếng Việt.
Kho riêng tư \texttt{pi-drkn-hinh} chỉ có một workflow và README.}
\end{frame}

%======================================================== 5. Thay đổi → triển khai
\begin{frame}{Một thay đổi đi từ máy tới người dùng}
\centering
\begin{tikzpicture}[x=1mm,y=1mm,node distance=4mm]
\node[box=MAC,text width=22mm] (e) {\textbf{Sửa mã}\\ nhánh riêng};
\node[box=MAC,text width=24mm,right=of e] (t) {\texttt{scripts/test-all.sh}\\ \textcolor{ink!60}{khoảng 2 phút, trên máy}};
\node[box=GH,text width=24mm,right=of t] (p) {\textbf{Pull request}\\ CI chạy lại kiểm thử};
\node[box=GH,text width=18mm,right=of p] (mg) {\textbf{Gộp} vào\\ \texttt{main}};
\node[box=MAC,text width=26mm,right=of mg] (d) {\texttt{scripts/deploy.sh}\\ kiểm thử lại → \texttt{clasp push}\\ → cùng bản triển khai};
\node[box=GG,text width=30mm,below=12mm of d] (g) {\textbf{Trên Google, tự động}\\ trong 1 giờ: \texttt{runSmoke}\\ (15 kịch bản) → thư nếu lỗi};
\node[box=GG,text width=30mm,left=8mm of g] (n) {\textbf{Mỗi đêm}: toàn bộ 110 kịch bản\\ trên Google → thư nếu lỗi};
\node[box=HT,text width=30mm,left=8mm of n] (u) {\textbf{Người dùng}\\ tải lại trang: bản mới\\ (địa chỉ không đổi)};
\draw[ar] (e)--(t); \draw[ar] (t)--(p); \draw[ar] (p)--(mg); \draw[ar] (mg)--(d); \draw[ar] (d)--(g); \draw[ar] (d.south west) to[out=-150,in=60] (u.north east);
\end{tikzpicture}

\vspace{3mm}\raggedright
\sm{\textbf{Ba lớp kiểm thử}: (1) trên máy và GitHub — chặn dữ liệu, 110 kịch bản trên Apps Script mô phỏng, trang thật trong Chromium với
mọi vai trò (kèm ảnh chụp thành hướng dẫn bằng hình \texttt{out/ui/index.html}); (2) trên Google — tự động như trên;
(3) bằng tay — chỉ những gì cần Google, hộp thư, máy Mac thật (\texttt{docs/testing.md}).
\texttt{deploy.sh} không đẩy mã nếu một kiểm thử hỏng.}
\end{frame}

%======================================================== 6. Triển khai từ đầu
\begin{frame}{Triển khai từ đầu (khoảng một buổi) — chi tiết: \texttt{docs/setup.md}}
\begin{columns}[T]
\begin{column}{0.5\textwidth}\sm{%
\textbf{Chuẩn bị}\\
\textbf{1.}\ Một tài khoản Google \textbf{của Pi} làm chủ (bật xác minh 2 bước). Mọi thứ dưới đây tạo bằng tài khoản này.\\
\textbf{2.}\ Máy có Node.js, \texttt{clasp} (\texttt{npm i -g @google/clasp}), \texttt{clasp login}; bật Apps Script API.
MacTeX nếu cần gói chế bản / dựng hình tay.\\
\textbf{3.}\ \texttt{git clone} kho mã; \texttt{npm install}; \texttt{scripts/test-all.sh} phải đạt.\\[2mm]
\textbf{Google}\\
\textbf{4.}\ \texttt{clasp create} hai dự án (\texttt{apps/main}, \texttt{apps/signin}).\\
\textbf{5.}\ \texttt{scripts/deploy.sh} → chép hai ID triển khai vào \texttt{deploy.local.env}.\\
\textbf{6.}\ Dự án chính: chạy \texttt{setup} → tạo Sheet, thư mục, in \texttt{SECRET}.\\
\textbf{7.}\ Dự án đăng nhập: \texttt{SECRET}, \texttt{B\_URL}; dự án chính: \texttt{SIGNIN\_URL}.}
\end{column}
\begin{column}{0.48\textwidth}\sm{%
\textbf{8.}\ Tab \texttt{Users}: email, vai trò của từng người.\\
\textbf{9.}\ \texttt{installReminderTrigger} (nhắc hạn).\\
\textbf{10.}\ Kiểm thử: \texttt{TEST\_USERS}, \texttt{setupTests}, \texttt{installAutoTests}.\\[2mm]
\textbf{GitHub (dựng hình tự động, tuỳ chọn)}\\
\textbf{11.}\ Kho riêng tư \texttt{pi-drkn-hinh} (chép workflow), mã thông báo fine-grained chỉ kho đó (Contents đọc/ghi)
→ \texttt{GITHUB\_REPO}, \texttt{GITHUB\_TOKEN}.\\[2mm]
\textbf{Dữ liệu}\\
\textbf{12.}\ Nhập hồ sơ cũ (\texttt{importLatest}) hoặc bắt đầu trống, thêm bài trên trang.\\
\textbf{13.}\ Gửi mọi người \textbf{địa chỉ /exec của ứng dụng đăng nhập}.\\[2mm]
\textcolor{warn!80!black}{Không bao giờ bấm „New deployment" trong trình soạn thảo: tạo địa chỉ mới, người dùng mất đường vào.}}
\end{column}
\end{columns}
\end{frame}

%======================================================== 7. Bàn giao
\begin{frame}{Nắm quyền mã nguồn và dữ liệu (bàn giao cho Pi)}
\begin{tikzpicture}[x=1mm,y=1mm]
\ccol{0}{GH}{1. Mã nguồn (GitHub)}{%
• Tạo \textbf{tổ chức GitHub của Pi}; chuyển hai kho vào đó: Settings → \emph{Transfer ownership}. Địa chỉ cũ tự chuyển hướng.\\
• Ít nhất \textbf{hai} quản trị viên, bật 2 bước.\\
• Kho \texttt{pi-drkn-hinh}: sửa \texttt{repository:} trong workflow sang tên mới; tạo \textbf{mã thông báo mới} của Pi, thay \texttt{GITHUB\_TOKEN}, \texttt{GITHUB\_REPO}.\\
• Mã đã công khai: ai cũng tải được; quyền ghi chỉ thành viên.}
\ccol{48}{GG}{2. Google (dự án, Sheet, Drive)}{%
• Chủ hiện tại chia sẻ cho tài khoản Pi rồi \emph{Transfer ownership}: hai dự án Apps Script, Sheet dữ liệu, Sheet kiểm thử, thư mục hình, chế bản.\\
• Tài khoản mới: \texttt{clasp login}, chạy \texttt{deploy.sh} (cùng ID triển khai) để ứng dụng chạy bằng quyền của mình; cho phép quyền.\\
• Trigger thuộc người tạo: chủ cũ xoá, chủ mới chạy lại \texttt{installReminderTrigger}, \texttt{installAutoTests}.}
\ccol{96}{warn}{3. Bí mật và thói quen}{%
• Đổi \texttt{SECRET} (cả hai dự án, cùng giá trị) — mọi phiên cũ hết hiệu lực.\\
• \texttt{deploy.local.env}, \texttt{.clasp.json} chỉ nằm trên máy người triển khai.\\
• Hai người biết triển khai; một người giữ quyền Quản trị dự phòng.\\
• Mã thông báo GitHub hết hạn tối đa 1 năm: ghi lịch đổi.\\
• Sao lưu: tự động mỗi đêm (30 bản, thư mục „Pi ĐRKN — sao lưu"); khôi phục: setup.md mục 15.}
\end{tikzpicture}
\end{frame}

%======================================================== 8. Vận hành
\begin{frame}{Vận hành: cái gì tự chạy, khi hỏng thì xem đâu}
\begin{columns}[T]
\begin{column}{0.5\textwidth}\sm{%
\textbf{Tự chạy (trigger của tài khoản chủ)}\\
• \texttt{sendReminders} — mỗi ngày 8 giờ: thư nhắc 3 và 1 ngày trước hạn.\\
• \texttt{autoTests} — mỗi giờ: bản mới → \texttt{runSmoke}; 2 giờ sáng → toàn bộ; kèm đồng bộ hình; 1 giờ sáng → sao lưu Sheet.\\
• \texttt{dongBoHinh} — sau khi lưu bài có TikZ, lặp mỗi 3 phút khi còn hình chờ GitHub.\\[2mm]
\textbf{Giới hạn của Google (tài khoản thường)}\\
• Một lần chạy tối đa 6 phút (kiểm thử, nhập tự chạy tiếp).\\
• Khoảng 100 người nhận thư mỗi ngày.\\
• Trigger tổng cộng 90 phút mỗi ngày.}
\end{column}
\begin{column}{0.48\textwidth}\sm{%
\textbf{Khi có sự cố}\\
• Người dùng không vào được → tab \texttt{Users} (vai trò, \texttt{hoat\_dong}), tab \texttt{Audit} (dòng „từ chối").\\
• Lỗi trên trang → trình soạn thảo Apps Script → \emph{Executions}.\\
• Thư kiểm thử báo lỗi → tab „Kết quả" của Sheet kiểm thử; chạy lại \texttt{runTests('<nhóm>')}.\\
• Hình không dựng → trang Hình (lỗi của LaTeX), kho \texttt{pi-drkn-hinh} → \emph{Actions}.\\
• Cần quay lại bản trước → \texttt{git checkout <commit>} rồi \texttt{deploy.sh}.\\[2mm]
\textbf{Chi phí}: không — Google và GitHub ở mức miễn phí.}
\end{column}
\end{columns}
\end{frame}

\end{document}
